% =============================================================================
% PREÁMBULO: Estadística con R — Fundamentos y Aplicaciones
% =============================================================================

% --- Codificación y lenguaje (LuaLaTeX maneja UTF-8 nativamente) -------------
\usepackage{fontspec}         % Gestión de fuentes OpenType con LuaLaTeX
\defaultfontfeatures{Ligatures=TeX}

% Babel v3.x (TeX Live 2025): carga idioma con \babelprovide
\usepackage{babel}
\babelprovide[main, import]{spanish}

% --- Tipografía --------------------------------------------------------------
\usepackage{lmodern}          % Latin Modern
\usepackage{microtype}        % Optimización tipográfica (kerning, protrusion)

% --- Matemáticas -------------------------------------------------------------
\usepackage{amsmath}          % Entornos equation, align, gather, etc.
\usepackage{amssymb}          % Símbolos matemáticos adicionales
\usepackage{amsthm}           % Entornos theorem, proof, definition
\usepackage{bm}               % Negritas en modo matemático (\bm)
\usepackage{mathtools}        % Extensiones de amsmath

% --- Definiciones de entornos matemáticos -----------------------------------
\theoremstyle{definition}
\newtheorem{definicion}{Definición}[chapter]
\newtheorem{ejemplo}{Ejemplo}[chapter]
\newtheorem{propiedad}{Propiedad}[chapter]

\theoremstyle{plain}
\newtheorem{teorema}{Teorema}[chapter]
\newtheorem{corolario}{Corolario}[chapter]
\newtheorem{lema}{Lema}[chapter]

\theoremstyle{remark}
\newtheorem{nota}{Nota}[chapter]
\newtheorem{observacion}{Observación}[chapter]

% --- Gráficos y figuras ------------------------------------------------------
\usepackage{graphicx}         % Inclusión de imágenes
\usepackage{float}            % Control de posición [H]
\usepackage{subcaption}       % Subfiguras
\graphicspath{{../output/}{./figuras/}}

% --- Tablas ------------------------------------------------------------------
\usepackage{booktabs}         % Líneas profesionales (\toprule, \midrule, \bottomrule)
\usepackage{array}            % Extensiones de columnas
\usepackage{longtable}        % Tablas que cruzan páginas (generadas por Pandoc)
\usepackage{multirow}         % Celdas que abarcan filas

% --- Código fuente R ---------------------------------------------------------
\usepackage{listings}
\usepackage{xcolor}

% Colores para el resaltado de sintaxis
\definecolor{codebg}{RGB}{248, 248, 248}
\definecolor{codeframe}{RGB}{200, 200, 200}
\definecolor{codecomment}{RGB}{106, 153, 85}
\definecolor{codekeyword}{RGB}{86, 156, 214}
\definecolor{codestring}{RGB}{206, 145, 120}
\definecolor{codenumber}{RGB}{181, 206, 168}
\definecolor{rblue}{RGB}{22, 92, 170}

% Estilo de código R
\lstdefinestyle{estilo_r}{
  language        = R,
  basicstyle      = \ttfamily\scriptsize,
  backgroundcolor = \color{codebg},
  frame           = single,
  framerule       = 0.5pt,
  rulecolor       = \color{codeframe},
  commentstyle    = \color{codecomment}\itshape,
  keywordstyle    = \color{codekeyword}\bfseries,
  stringstyle     = \color{codestring},
  numberstyle     = \tiny\color{codenumber},
  numbers         = left,
  numbersep       = 8pt,
  stepnumber      = 1,
  tabsize         = 2,
  showstringspaces= false,
  breaklines      = true,
  breakatwhitespace= true,
  captionpos      = b,
  morekeywords    = {TRUE, FALSE, NULL, NA, NaN, Inf,
                     library, require, install.packages,
                     function, return, if, else, for, while, repeat,
                     in, next, break,
                     data.frame, list, vector, matrix, array, factor,
                     read.csv, write.csv, source, setwd, getwd,
                     ggplot, aes, geom_point, geom_line, geom_bar,
                     geom_histogram, geom_boxplot, geom_smooth,
                     labs, theme_minimal, ggsave, scale_color_viridis_d,
                     lm, glm, aov, summary, coef, residuals, fitted,
                     t.test, wilcox.test, shapiro.test, cor, var.test,
                     mean, median, sd, var, IQR, quantile, range,
                     table, prop.table, tapply, sapply, lapply, apply,
                     filter, select, mutate, group_by, summarise, arrange,
                     pivot_wider, pivot_longer, left_join, inner_join,
                     set.seed, runif, rnorm, sample, rep, seq},
  literate        =
    {á}{{\'a}}1 {é}{{\'e}}1 {í}{{\'{\i}}}1 {ó}{{\'o}}1 {ú}{{\'u}}1
    {Á}{{\'A}}1 {É}{{\'E}}1 {Í}{{\'I}}1 {Ó}{{\'O}}1 {Ú}{{\'U}}1
    {ñ}{{\~n}}1 {Ñ}{{\~N}}1 {ü}{{\"u}}1 {Ü}{{\"U}}1
}

\lstset{style=estilo_r}

% Entorno para código inline
\newcommand{\rcode}[1]{\texttt{\small\color{rblue}#1}}

% --- Hiperenlaces y metadatos PDF -------------------------------------------
\usepackage{hyperref}
\hypersetup{
  pdftitle    = {Estadística con R: Fundamentos y Aplicaciones},
  pdfauthor   = {Proyecto Libro Caribe Colombiano},
  pdfsubject  = {Estadística, R, Análisis de Datos},
  pdfkeywords = {estadística descriptiva, inferencia, regresión, R, tidyverse,
                 Caribe colombiano, Sierra Nevada, Cesar, Barranquilla},
  colorlinks  = true,
  linkcolor   = rblue,
  urlcolor    = rblue,
  citecolor   = rblue,
  bookmarks   = true,
  bookmarksnumbered = true
}

% --- Diseño de página --------------------------------------------------------
\usepackage{geometry}
\geometry{
  letterpaper,                 % Formato carta (21.59 × 27.94 cm)
  top        = 2.5cm,
  bottom     = 2.5cm,
  left       = 3.0cm,          % Margen izquierdo / interior (encuadernación)
  right      = 2.5cm,          % Margen derecho / exterior
  headheight = 14pt
}

% --- Encabezados y pies de página -------------------------------------------
\usepackage{fancyhdr}
\pagestyle{fancy}
\fancyhf{}
\fancyhead[LE]{\small\thepage \quad \textit{Estadística con R}}
\fancyhead[RO]{\small\textit{\nouppercase{\rightmark}} \quad \thepage}
\fancyfoot[C]{}
\renewcommand{\headrulewidth}{0.4pt}

% --- Espaciado y párrafos ----------------------------------------------------
\usepackage{setspace}
\setstretch{1.15}              % Interlineado compacto (1.15x, vs 1.5x anterior)
\setlength{\parskip}{4pt}
\setlength{\parindent}{0pt}   % Sin sangría de primera línea (estilo moderno)

% --- Cajas de texto (notas, advertencias) ------------------------------------
\usepackage{tcolorbox}
\tcbuselibrary{skins, breakable}

\newtcolorbox{notabox}[1][]{
  enhanced, breakable,
  colback  = blue!5!white,
  colframe = rblue,
  title    = {#1},
  fonttitle= \bfseries\small,
  left=6pt, right=6pt, top=4pt, bottom=4pt
}

\newtcolorbox{codigobox}[1][]{
  enhanced, breakable,
  colback  = codebg,
  colframe = codeframe,
  title    = {#1},
  fonttitle= \ttfamily\small\color{rblue},
  left=4pt, right=4pt, top=2pt, bottom=2pt
}

% --- Referencias bibliográficas ----------------------------------------------
\usepackage[authoryear, round]{natbib}
\bibliographystyle{apalike}

% --- Compatibilidad con Pandoc -----------------------------------------------
% \tightlist: listas compactas generadas por Pandoc
\providecommand{\tightlist}{%
  \setlength{\itemsep}{2pt}\setlength{\parskip}{0pt}}
% \passthrough: código inline (texttt es seguro dentro de tablas)
\providecommand{\passthrough}[1]{\texttt{#1}}

% --- Comandos personalizados -------------------------------------------------
\newcommand{\E}[1]{\mathbb{E}\left[#1\right]}          % Esperanza
\newcommand{\Var}[1]{\mathrm{Var}\left(#1\right)}      % Varianza
\newcommand{\Cov}[2]{\mathrm{Cov}\left(#1,#2\right)}   % Covarianza
\newcommand{\N}[2]{\mathcal{N}\left(#2,#2\right)}      % Normal
\newcommand{\iid}{\overset{\text{iid}}{\sim}}           % iid
\newcommand{\pvalue}{p\text{-valor}}                    % p-valor
\newcommand{\dataset}[1]{\texttt{#1}}                   % Nombre de dataset
